% ============================================================
%  Kapitel 4: MadoLudus -- In-Window-Slideshow (v0.5.42)
% ============================================================
\chapter{Tab 3 -- MadoLudus: Konfiguration und Wiedergabefenster}

% ============================================================
\section{Funktionsweise}
% ============================================================

MadoLudus ist eine Diashow, die in einem \emph{eigenen,
rahmenlosen Fenster} (\klasse{MadoludusWindow}) laeuft.
Der Tab selbst ist nur die Konfigurationsoberflaeche -- das
Fenster oeffnet sich auf Knopfdruck.

\begin{neuin}[title={v0.5.42: Kompletter Refactor}]
Der fruehre \klasse{MadolodosTab} war ein Monolith:
Konfiguration und Wiedergabe in einer Klasse.
Ab v0.5.42 sind zwei Klassen getrennt:
\begin{itemize}[noitemsep]
  \item \klasse{MadoludusTab} -- Konfigurations-UI
    im Hauptfenster (bleibt immer sichtbar)
  \item \klasse{MadoludusWindow} -- Eigenstaendiges
    \klasse{QMainWindow}, rahmenlos, immer im Vordergrund,
    per Maus verschiebbar
\end{itemize}
Neu ebenfalls: Shuffle-Button (\textbf{R}-Taste) und
FF-Rate-Auswahl (1/sec, 1/2 sec, 1/N Frames) als Combobox.
\end{neuin}

\textbf{Betriebsmodi fuer Videos:}
\begin{center}
\small
\begin{tabular}{lll}
\toprule
\textbf{Modus} & \textbf{Bedingung} & \textbf{Mechanismus} \\
\midrule
Native Wiedergabe & FF-Mode aus, Multimedia vorhanden &
  \klasse{QMediaPlayer} + \klasse{QVideoWidget} \\
FF/Flip-Book & FF-Mode an (Standard) &
  ffmpeg-Frames via \klasse{VideoFramesWorker} \\
Fallback & kein Multimedia (\makro{TV\_NO\_MULTIMEDIA}) &
  immer ffmpeg-Frames \\
\bottomrule
\end{tabular}
\end{center}

\textbf{Tastaturkuerzel in MadoludusWindow:}
\begin{center}
\small
\begin{tabular}{ll}
\toprule
\textbf{Taste} & \textbf{Aktion} \\
\midrule
Space         & Play / Pause \\
Pfeil rechts  & Naechste Datei \\
Pfeil links   & Vorherige Datei (Verlauf-bewusst) \\
Shift+Rechts  & Video: +5 Sekunden (nur native) \\
Shift+Links   & Video: --5 Sekunden (nur native) \\
M             & Stummschaltung umschalten \\
F             & FF-Modus (Secret Mode zum Deaktivieren) \\
R             & Zufaellige / sequenzielle Reihenfolge \\
Esc           & Fenster schliessen \\
\bottomrule
\end{tabular}
\end{center}

% ============================================================
\section{File by File}
% ============================================================

\subsection{madoluduswindow.h -- MadoludusConfig-Struct}

\begin{lstlisting}[style=header,
  caption={madoluduswindow.h -- MadoludusConfig und Klassendeklaration}]
// Konfiguration, die MadoludusTab an MadoludusWindow uebergibt
struct MadoludusConfig {
    bool   showImages   = true;
    bool   showVideos   = true;
    bool   ffMode       = true;
    // Wert fuer Frame-Extraktion: Sekunden oder Anzahl
    double ffValue      = 10.0;
    // false = IntervalSeconds, true = FixedCount
    bool   ffFixedCount = false;
    bool   muted        = false;
    bool   randomMode   = false;
    bool   secretMode   = false;
};

class MadoludusWindow : public QMainWindow
{
    Q_OBJECT
public:
    explicit MadoludusWindow(
        const QStringList&    mediaFiles,
        const MadoludusConfig& cfg,
        QWidget*              parent = nullptr);
    ~MadoludusWindow() override;

signals:
    void windowClosed();
    void currentFileChanged(const QString& path);

protected:
    void keyPressEvent(QKeyEvent*  event) override;
    void resizeEvent(QResizeEvent* event) override;
    void mousePressEvent(QMouseEvent* event) override;
    void mouseMoveEvent(QMouseEvent*  event) override;
    void closeEvent(QCloseEvent*   event) override;

private slots:
    void onNextFile();
    void onPrevFile();
    void onPlayPause();
    void onMuteToggle();
    void onFFToggle();
    void onShuffleToggle();    // v0.5.42
    void onAutoAdvance();
    void onVideoEnded();
    void onFFFramesReady(bool ok, QStringList framePaths);
    void cycleFFFrame();

private:
    void buildUi();
    void loadFile(int index);
    void loadImage(const QString& path);
    void loadVideo(const QString& path);
    int  nextIndex() const;
    void startAutoAdvance();
    void stopAutoAdvance();
    void startFFExtraction(const QString& path);
    void stopFFExtraction();
    void updateButtonStates();
    void scaleCurrentImage();

    QStringList     m_mediaFiles;
    MadoludusConfig m_cfg;
    int             m_index        = -1;
    bool            m_paused       = false;
    bool            m_isVideo      = false;
    bool            m_usingFlipbook = false;

    // Navigations-Verlauf (Previous funktioniert auch im Zufallsmodus)
    QList<int> m_order;
    int        m_orderPos = -1;

    QWidget*     m_central   = nullptr;
    QLabel*      m_display   = nullptr;  // Bild / Flip-Book
    QLabel*      m_lblPath   = nullptr;  // Pfad unten
    MadoOverlay* m_overlay   = nullptr;  // linke Steuerleiste

    QTimer* m_advTimer     = nullptr;   // Auto-Advance
    QTimer* m_ffCycleTimer = nullptr;   // Frame-Wechsel

    QThread*    m_ffThread   = nullptr;
    QStringList m_ffFrames;
    int         m_ffFrameIdx = 0;
    QString     m_ffTmpDir;

#ifndef TV_NO_MULTIMEDIA
    QWidget*      m_videoContainer = nullptr;
    QMediaPlayer* m_player         = nullptr;
    QAudioOutput* m_audio          = nullptr;
    QVideoWidget* m_videoW         = nullptr;
#endif

    // Drag-to-move (rahmenlos)
    QPoint m_dragStart;
    bool   m_dragging = false;

    static constexpr int kImageDelay   = 5000; // ms
    static constexpr int kImageDelayFF = 2000;
    static constexpr int kFFCycleMs    = 600;
    static constexpr int kSeekStep     = 5000;
};
\end{lstlisting}

\subsubsection{C++-Analyse: Struct mit Default-Werten}

\begin{lstlisting}
struct MadoludusConfig {
    bool   ffMode  = true;
    double ffValue = 10.0;
    // ...
};
\end{lstlisting}

Ein \verb|struct| in C++ ist wie eine \verb|class|, aber
mit Standard-Sichtbarkeit \verb|public|. Die Member
werden direkt initialisiert (C++11). Das bedeutet:
\begin{lstlisting}
MadoludusConfig cfg;   // alle Member auf Default-Werte
cfg.muted = true;      // einzelne anpassen
\end{lstlisting}
Dieses Muster heisst \emph{Aggregate Initialization} --
kein Konstruktor noetig, solange kein
Verhalten gekapselt werden muss.

\subsubsection{C++-Analyse: Navigations-Verlauf}

\begin{lstlisting}[caption={nextIndex und Verlaufslogik}]
int MadoludusWindow::nextIndex() const {
    if (m_mediaFiles.size() <= 1)
        return qMax(0, m_index);
    if (!m_cfg.randomMode)
        return (m_index + 1) % m_mediaFiles.size();

    // Zufaelligen Index waehlen (nicht den aktuellen)
    std::mt19937 rng(std::random_device{}());
    std::uniform_int_distribution<int>
        dist(0, m_mediaFiles.size() - 1);
    int n;
    do { n = dist(rng); } while (n == m_index);
    return n;
}

void MadoludusWindow::onNextFile() {
    if (m_mediaFiles.isEmpty()) return;
    int next;
    if (m_orderPos < m_order.size() - 1) {
        // Verlauf vorwaerts (Nutzer hatte zurueck
        // navigiert und geht nun wieder vor)
        ++m_orderPos;
        next = m_order.at(m_orderPos);
    } else {
        // Neuen Index berechnen und in Verlauf aufnehmen
        next = nextIndex();
        m_order.append(next);
        m_orderPos = m_order.size() - 1;
    }
    loadFile(next);
}
\end{lstlisting}

Das Verlaufsystem mit \member{m\_order} (Liste der
angezeigten Indizes) und \member{m\_orderPos} (Position
darin) erlaubt beliebiges Vor- und Rueckwaertsnavigieren
-- auch im Zufallsmodus wiederholt man exakt die vorher
gezeigten Dateien.

\subsubsection{Qt-Analyse: Rahmenloses Fenster und Drag-to-Move}

\begin{lstlisting}[caption={Konstruktor: Fenster-Flags}]
setWindowFlags(Qt::Window
             | Qt::FramelessWindowHint
             | Qt::WindowStaysOnTopHint);
setAttribute(Qt::WA_DeleteOnClose);
setStyleSheet("background: black;");
\end{lstlisting}

\begin{itemize}[noitemsep]
  \item \verb|FramelessWindowHint| -- kein Titelbalken,
    kein OS-Rahmen; Fenster muss manuell bewegbar gemacht
    werden
  \item \verb|WindowStaysOnTopHint| -- immer im Vordergrund
    (auch vor anderen Anwendungen)
  \item \verb|WA_DeleteOnClose| -- das Objekt wird beim
    Schliessen automatisch geloescht (kein Speicherleck)
\end{itemize}

Drag-to-Move-Implementierung:
\begin{lstlisting}[caption={Mouse-Events fuer Drag-to-Move}]
void MadoludusWindow::mousePressEvent(QMouseEvent* event) {
    if (event->button() == Qt::LeftButton &&
        // Nicht dragging wenn im Overlay-Bereich
        !m_overlay->geometry().contains(event->pos()))
    {
        m_dragging  = true;
        m_dragStart = event->globalPosition().toPoint()
                      - frameGeometry().topLeft();
    }
    QMainWindow::mousePressEvent(event);
}

void MadoludusWindow::mouseMoveEvent(QMouseEvent* event) {
    if (m_dragging && (event->buttons() & Qt::LeftButton))
        move(event->globalPosition().toPoint() - m_dragStart);
    QMainWindow::mouseMoveEvent(event);
}
\end{lstlisting}

\verb|globalPosition()| liefert die absolute Bildschirmposition
der Maus (Qt6: gibt \verb|QPointF| zurueck, daher
\verb|.toPoint()|). Der Offset \member{m\_dragStart}
ist die Differenz zwischen Mausposition und Fenster-Ecke --
so bewegt sich das Fenster nicht sprunghaft.

\subsubsection{Qt-Analyse: FF-Flip-Book-Pipeline}

\begin{lstlisting}[caption={startFFExtraction: VideoFramesWorker starten}]
void MadoludusWindow::startFFExtraction(const QString& path)
{
    stopFFExtraction();

    const VideoFramesWorker::Mode mode =
        m_cfg.ffFixedCount
        ? VideoFramesWorker::Mode::FixedCount
        : VideoFramesWorker::Mode::IntervalSeconds;

    auto* worker = new VideoFramesWorker(
        path, mode, m_cfg.ffValue);
    m_ffTmpDir = worker->tmpDir();
    auto* thread = new QThread;
    m_ffThread = thread;
    worker->moveToThread(thread);

    connect(thread, &QThread::started,
            worker, &VideoFramesWorker::run);
    connect(worker, &VideoFramesWorker::resultReady,
            this,   &MadoludusWindow::onFFFramesReady);
    connect(worker, &VideoFramesWorker::resultReady,
            thread, &QThread::quit);
    connect(worker, &VideoFramesWorker::resultReady,
            worker, &QObject::deleteLater);
    connect(thread, &QThread::finished,
            thread, &QObject::deleteLater);
    thread->start();
}

void MadoludusWindow::onFFFramesReady(
    bool ok, QStringList framePaths)
{
    m_ffThread = nullptr;
    if (!ok || framePaths.isEmpty()) {
        m_display->setText(
            "(kein Video-Preview -- ffmpeg installieren)");
        startAutoAdvance();
        return;
    }
    m_ffFrames    = framePaths;
    m_ffFrameIdx  = 0;
    cycleFFFrame();
    if (!m_paused) {
        m_ffCycleTimer->start(kFFCycleMs);
        // Auto-Advance nach allen Frames einmal gesehen
        m_advTimer->start(
            qMax(kImageDelayFF,
                 kFFCycleMs * m_ffFrames.size()));
    }
}
\end{lstlisting}

\subsubsection{Qt-Analyse: scaleCurrentImage beim Resize}

\begin{lstlisting}[caption={resizeEvent: Bild neu skalieren}]
void MadoludusWindow::resizeEvent(QResizeEvent* event) {
    QMainWindow::resizeEvent(event);
    scaleCurrentImage();
}

void MadoludusWindow::scaleCurrentImage() {
    if (!m_display->isVisible() ||
        m_display->pixmap().isNull()) return;

    if (m_usingFlipbook && !m_ffFrames.isEmpty()) {
        // Letzten angezeigten Frame neu skalieren
        const int shown =
            (m_ffFrameIdx + m_ffFrames.size() - 1)
            % m_ffFrames.size();
        const QPixmap pm(m_ffFrames.at(shown));
        if (!pm.isNull())
            m_display->setPixmap(pm.scaled(
                m_display->size(),
                Qt::KeepAspectRatio,
                Qt::SmoothTransformation));
    } else if (!m_isVideo && m_index >= 0) {
        // Statisches Bild neu laden und skalieren
        QImageReader reader(m_mediaFiles.at(m_index));
        const QSize win = m_display->size();
        if (!win.isEmpty()) {
            reader.setScaledSize(win);
            const QImage img = reader.read();
            if (!img.isNull())
                m_display->setPixmap(
                    QPixmap::fromImage(img).scaled(
                        win, Qt::KeepAspectRatio,
                        Qt::SmoothTransformation));
        }
    }
}
\end{lstlisting}

% ----------------------------------------------------------
\subsection{madoludustab.h / madoludustab.cpp}
% ----------------------------------------------------------

\subsubsection{Klassendeklaration}

\begin{lstlisting}[style=header, caption={madoludustab.h -- Kern}]
class MadoludusTab : public QWidget {
    Q_OBJECT
public:
    explicit MadoludusTab(
        std::function<QString()> getGlobalDir,
        QWidget* parent = nullptr);
    ~MadoludusTab() override;  // stoppt ggf. Scan-Thread

    void onDirectoryChanged(const QString& path,
                            const QStringList& allFiles = {});
    void setSecretMode(bool on);

private slots:
    void showPreview(const QString& path);
    void onStartClicked();
    void onScanDone(bool ok,
                    QStringList imageFiles,
                    QStringList allFiles);
    void onWindowClosed();
    void openWindow(const QStringList& files);

private:
    void buildUi();
    void scanDir(const QString& path);
    MadoludusConfig buildConfig() const;
    QStringList filteredFiles() const;

    std::function<QString()> m_getGlobalDir;
    QStringList m_allFiles;
    bool        m_secretMode = false;

    QCheckBox*   m_chkImages  = nullptr;
    QCheckBox*   m_chkVideos  = nullptr;
    QCheckBox*   m_chkFF      = nullptr;
    QComboBox*   m_cmbFFRate  = nullptr;  // v0.5.42
    QCheckBox*   m_chkMuted   = nullptr;
    QCheckBox*   m_chkRandom  = nullptr;  // v0.5.42
    QPushButton* m_btnPreview = nullptr;
    QPushButton* m_btnStart   = nullptr;
    QLabel*      m_lblStatus  = nullptr;
    QProgressBar* m_pb        = nullptr;

    VideoPreviewWidget* m_preview = nullptr;
    QLabel*             m_lblType = nullptr;
    QThread*        m_scanThread  = nullptr;
    MadoludusWindow* m_window     = nullptr;
};
\end{lstlisting}

\subsubsection{C++-Analyse: buildConfig und FF-Rate-Tabelle}

\begin{lstlisting}[caption={FF-Rate-Tabelle und buildConfig}]
// Tabelle aller FF-Raten-Optionen:
struct FFRateEntry {
    const char* label;
    bool   fixedCount; // false=Intervall, true=Anzahl
    double value;
};
static const FFRateEntry kFFRates[] = {
    { "1/sec",       false, 1.0  },
    { "1/2 sec",     false, 2.0  },
    { "1/5 sec",     false, 5.0  },
    { "1/10 sec",    false, 10.0 },
    { "1/2 frames",  true,  2.0  },
    { "1/10 frames", true,  10.0 },
    { "1/30 frames", true,  30.0 },
    { "1/60 frames", true,  60.0 },
};
static constexpr int kFFRateCount =
    static_cast<int>(sizeof(kFFRates)/sizeof(kFFRates[0]));

// Config aus UI-Zustand bauen:
MadoludusConfig MadoludusTab::buildConfig() const {
    MadoludusConfig cfg;
    cfg.showImages  = m_chkImages->isChecked();
    cfg.showVideos  = m_chkVideos->isChecked();
    cfg.ffMode      = m_chkFF->isChecked();
    cfg.muted       = m_chkMuted->isChecked();
    cfg.randomMode  = m_chkRandom->isChecked();

    const int idx = m_cmbFFRate->currentIndex();
    if (idx >= 0 && idx < kFFRateCount) {
        cfg.ffFixedCount = kFFRates[idx].fixedCount;
        cfg.ffValue      = kFFRates[idx].value;
    } else {
        cfg.ffFixedCount = false;
        cfg.ffValue      = 10.0;
    }
    return cfg;
}
\end{lstlisting}

\verb|sizeof(kFFRates)/sizeof(kFFRates[0])| berechnet
die Arraylaenge zur Compile-Zeit ohne hartcodierten Wert --
bei jedem neuen Eintrag in der Tabelle stimmt der Count
automatisch.

\subsubsection{Qt-Analyse: openWindow und Signal-Rueckkanal}

\begin{lstlisting}[caption={openWindow: Fenster starten und verbinden}]
void MadoludusTab::openWindow(const QStringList& files) {
    MadoludusConfig cfg = buildConfig();
    cfg.secretMode = m_secretMode;
    m_window = new MadoludusWindow(files, cfg);

    // Startgroesse: halber Bildschirm, linke obere Ecke
    const QRect scr =
        QGuiApplication::primaryScreen()
        ->availableGeometry();
    m_window->setGeometry(
        0, 0, scr.width()/2, scr.height()/2);

    // Rueckkanal: Fenster meldet geschlossen
    connect(m_window, &MadoludusWindow::windowClosed,
            this,     &MadoludusTab::onWindowClosed);

    // Rueckkanal: Datei geaendert -> Preview aktualisieren
    connect(m_window,
            &MadoludusWindow::currentFileChanged,
            this,
            [this](const QString& path) {
                showPreview(path);
            });

    m_window->show();
    m_lblStatus->setText(
        QString("MadoLudus running -- %1 file(s).")
        .arg(files.size()));
}

void MadoludusTab::onWindowClosed() {
    m_window = nullptr;  // Zeiger freigeben (WA_DeleteOnClose)
    m_lblStatus->setText("MadoLudus closed.");
}
\end{lstlisting}

\verb|WA_DeleteOnClose| (gesetzt im \klasse{MadoludusWindow}-Konstruktor)
loescht das Objekt beim Schliessen automatisch.
Der Zeiger \member{m\_window} wird im \slot{onWindowClosed}-Slot
auf \verb|nullptr| gesetzt -- sonst haette man einen
Dangling Pointer.

% ----------------------------------------------------------
\subsection{madooverlay.h / madooverlay.cpp}
% ----------------------------------------------------------

\klasse{MadoOverlay} ist die halbtransparente Steuerleiste
im \klasse{MadoludusWindow}. Sie ist ein eigenstaendiges
\klasse{QWidget}, das als Kind im Fenster positioniert wird.

\begin{lstlisting}[style=header, caption={madooverlay.h}]
class MadoOverlay : public QWidget {
    Q_OBJECT
public:
    explicit MadoOverlay(QWidget* parent = nullptr);

    void setPaused(bool paused);   // Play/Pause-Icon
    void setMuted(bool muted);     // Lautsprecher-Icon
    void setFFMode(bool ff);       // FF-Button-Farbe
    void setRandomMode(bool on);   // v0.5.42: Shuffle-Icon
    void setFileInfo(int index, int total,
                     const QString& name);
    void setSecretMode(bool on);

signals:
    void playPauseClicked();
    void prevClicked();
    void nextClicked();
    void muteClicked();
    void ffClicked();
    void shuffleClicked();         // v0.5.42

protected:
    void paintEvent(QPaintEvent* event) override;

private:
    QPushButton* m_btnPrev    = nullptr;
    QPushButton* m_btnPlay    = nullptr;
    QPushButton* m_btnNext    = nullptr;
    QPushButton* m_btnMute    = nullptr;
    QPushButton* m_btnFF      = nullptr;
    QPushButton* m_btnShuffle = nullptr;  // v0.5.42
    QLabel*      m_lblInfo    = nullptr;
};
\end{lstlisting}

\subsubsection{Qt-Analyse: Custom paintEvent}

\begin{lstlisting}[caption={MadoOverlay::paintEvent -- dunkle abgerundete Pille}]
void MadoOverlay::paintEvent(QPaintEvent*) {
    QPainter p(this);
    p.setRenderHint(QPainter::Antialiasing);
    // Schwarz mit 63% Deckkraft (Alpha 160/255)
    p.setBrush(QColor(0, 0, 0, 160));
    p.setPen(Qt::NoPen);
    p.drawRoundedRect(
        rect().adjusted(4, 4, -4, -4), 10, 10);
}
\end{lstlisting}

\verb|QColor(r, g, b, alpha)| -- der vierte Parameter
ist der Alpha-Kanal (0 = voellig transparent,
255 = voellig opak). \verb|drawRoundedRect| mit Radius 10
erzeugt die "Pille"-Form. \verb|Qt::NoPen| verhindert
einen Rahmen.

\begin{hinweis}
\textbf{Wann paintEvent ueberschreiben?}
Normalerweise reicht es, ein \klasse{QWidget} per
\verb|setStyleSheet()| zu stylen.
\verb|paintEvent| ueberschreibt man nur, wenn man
volle Kontrolle ueber das Zeichnen braucht -- hier
fuer die halbtransparente Pille, die kein
Stylesheet erreichen kann.
\end{hinweis}

\subsubsection{Qt-Analyse: Unicode-Icons in Buttons}

\begin{lstlisting}[caption={makeBtn: Emoji-Symbole als Button-Texte}]
static QPushButton* makeBtn(const QString& text,
                             QWidget* parent) {
    auto* b = new QPushButton(text, parent);
    b->setFixedSize(40, 40);
    b->setFlat(true);
    b->setStyleSheet(
        "QPushButton {"
        "  font-size: 18px;"
        "  color: white;"
        "  background: transparent;"
        "  border: none;"
        "  border-radius: 6px;"
        "}"
        "QPushButton:hover {"
        "  background: rgba(255,255,255,40);"
        "}");
    return b;
}
// Aufruf:
m_btnPlay    = makeBtn("\u23f8", this);  // Pause-Symbol
m_btnMute    = makeBtn("\U0001f507", this); // Stummschalten
m_btnShuffle = makeBtn("\U0001f500", this); // Shuffle
\end{lstlisting}

Qt stellt Unicode-Zeichen als String-Literale dar.
\verb|\uXXXX| (4 Hex-Ziffern) fuer BMP-Zeichen,
\verb|\UXXXXXXXX| (8 Hex-Ziffern) fuer Zeichen
ausserhalb der BMP (Supplementary Planes, z.B. Emojis).
Die Schriftart muss die Zeichen unterstuetzen --
auf Windows liefert "Segoe UI Emoji" alle gaengigen Emojis.
